\subsection{From Method Calls to Message Protocols}

In a familiar OOP language, a call such as \code{account.balance()}
selects a method, executes its body, and returns a value to the caller.
ABCL/R2 separates these jobs. A sender constructs a data value, sends it
to an object, and may arrange a destination for a later reply.
The receiving object matches that data against its scripts.
A \emph{script} is a handler: a message pattern followed by expressions
to execute when the pattern matches.

\begin{longtable}{L{0.28\textwidth}L{0.62\textwidth}}
\textbf{Familiar idea} & \textbf{ABCL/R2 counterpart and difference}\\\hline
Object fields & State bindings belonging to an object; outer lexical
bindings can also be shared.\\
Method selection & First matching message pattern, not necessarily a
method name plus a fixed argument signature.\\
Method arguments & Data inside the message; pattern variables extract
the parts the handler needs.\\
Return value & An explicit reply message, not the handler's final value.\\
Inheritance & Can be implemented by a reflective protocol; the original
bird/emu example does not declare a superclass.\\
\end{longtable}

A \emph{protocol} specifies the messages an object accepts and what its
clients may expect: replies, state changes, or further messages.
For an account, \code{[:deposit 10]} and \code{:balance} can be two
protocol messages. The colon is part of a Lisp keyword, not special
method-call punctuation; the next chapter explains keywords.

\subsection{An Actor-Like Mental Model}

An \emph{actor} is a useful introductory model of an independently
addressable participant that reacts to messages and can send more messages.
For an ordinary ABCL/R2 object, imagine an inbox and a handler processing
one request at a time. Other objects have their own work.
\emph{Concurrent} means their activities can be outstanding together;
it does not require simultaneous execution on different CPUs.

The distinction becomes concrete in this sequence:
\begin{enumerate}[nosep]
  \item A client sends a deposit request to an account.
  \item The request enters the ordinary receiving object's message machinery.
  \item The account selects a matching handler and updates its balance.
  \item If the protocol requires it, the handler explicitly sends a reply.
  \item The client uses that reply, or continues without waiting if it sent
        asynchronously.
\end{enumerate}
Sending is therefore not the same event as finishing the requested work.
An asynchronous sender does not automatically receive a completion notice.
Conversely, a waiting sender cannot finish its request unless someone
sends the required reply.

This is an actor-like explanation, not a claim that ABCL/R2 enforces pure
actor isolation. Its lexical bindings can be shared between objects.
\emph{Reflection} means that a program can inspect or control parts of
its own execution machinery. It can therefore change the ordinary
selection and execution of handlers.
Learn the ordinary model first; later chapters state where it is extended.

\subsection{Three Questions to Ask While Reading}

For every send in an original program, ask:
\begin{enumerate}[nosep]
  \item Which object receives this message, and which pattern matches it?
  \item Does the sender wait, and where will any reply be sent?
  \item What establishes completion: one reply, a count, or a terminal message?
\end{enumerate}
In \filepath{examples/misc/future-ex.abclr2}, these questions distinguish
the tree walker, the object consuming its replies, and the printer.
In \filepath{examples/misc/inherit.abclr2}, they distinguish an application
object from the metaobject running its dispatch protocol.
You do not need to understand those files yet; these are our first
original-source reading goals.

\paragraph{A route through the book.}
Work in stages rather than trying to absorb every library at once:
\begin{enumerate}[nosep]
  \item \emph{Build the ordinary model.} Follow
        Sections~\ref{sec:first-session}--\ref{sec:files}, trying the
        small sessions and exercises. At the bank account
        (Section~\ref{sec:account}), explain which pattern handles each
        request and why the caller receives a reply.
  \item \emph{Follow results over time.} Study futures in
        Section~\ref{sec:futures}, then trace the original tree walker
        in Section~\ref{sec:original-future}. Explain why its terminal
        message, not an empty buffer or a printed line, marks completion.
        Express messages are an advanced extension; you may return to
        that experiment after learning reflection.
  \item \emph{Open the execution machinery.} Read
        Section~\ref{sec:reflection}, then the original lazy-control,
        inheritance, and interpreter projects in
        Section~\ref{sec:originals}. Keep the application reply separate
        from the interpreter's execution-completion message.
  \item \emph{Combine the mechanisms.} Work through the library and
        debugging projects, then Pcontrol and Time Warp
        (Sections~\ref{sec:pcontrol} and \ref{sec:timewarp}).
        These are advanced reading projects, not prerequisites for
        ordinary object programming. Return to the file map in
        Appendix~\ref{sec:source-map} when following their helper objects.
\end{enumerate}
The job-service project (Section~\ref{sec:project}) is a parallel practice
track: it combines ordinary objects and reply protocols, so you can try
it after the foundations without first mastering the reflective projects.
Use the worked exercise answers to check your reasoning, not just the
final values. At each stage, being able to explain a message's recipient,
reply destination, and completion condition is the checkpoint for moving on.
